% Part: first-order-logic
% Chapter: models-theories
% Section: set-theory

\documentclass[../../../include/open-logic-section]{subfiles}

\begin{document}

\olfileid{fol}{mat}{set}

\olsection{Teori Himpunan}

Meh kabeh matematika bisa dikembangake ing teori himpunan.
Ngembangake matematika ing teori iki nyakup sawetara prakara.
Kapisan, dibutuhake himpunan aksioma kanggo relasi~$\in$.
Wis dikembangake maneka sistem aksioma, kang kadhang
nduweni sipat-sipat~$\in$ kang padha nentang.
Sistem aksioma kang dikenal minangka $\Log{ZFC}$,
teori himpunan Zermelo--Fraenkel kanthi aksioma pilihan,
ngadeg dhewe: sistem iki kang paling akeh digunakake
lan ditliti, amarga aksioma-aksiomane kabukten cukup
kanggo mbuktekake meh kabeh bab kang dikarepake
para matematikawan bisa dibuktekake. Nanging sadurunge
iku bisa ditetepake, luwih dhisik kudu dijlentrehake
kepriye kita bisa \emph{nyatakake} kabeh bab kang
kepengin dinyatakake dening para matematikawan.
Kaping pisan, basa iki ora ngemot !!{constant}s
utawa !!{function}s, mula sepisanan katon ora cetha
kepriye kita bisa ngrembug himpunan tartamtu
(kayata $\emptyset$ utawa $\Nat$), operasi tumrap
himpunan (kayata $X \cup Y$ lan $\Pow{X}$),
apa maneh konstruksi liyane kang nyakup barang
liya saliyane himpunan, kayata relasi lan fungsi.

Kanggo miwiti, "dadi anggota saka" dudu siji-sijine
relasi kang kita gatekake: "dadi subhimpunan saka"
katon meh padha wigatine. Nanging kita bisa
\emph{nemtokake} "dadi subhimpunan saka" lumantar
"dadi anggota saka". Kanggo nindakake iki,
kita kudu nemokake !!a{formula}~$!A(x, y)$
ing basa teori himpunan kang dicukupi dening
pasangan himpunan~$\tuple{X, Y}$ yen lan mung
yen $X \subseteq Y$. Dene $X$ iku subhimpunan
saka $Y$ yen lan mung yen kabeh !!{element}s
saka~$X$ uga !!{element}s saka~$Y$. Mula,
kita bisa nemtokake $\subseteq$ kanthi formula
\[
\lforall[z][(z \in x \lif z \in y)]
\]
Saiki, kapan wae kita arep nggunakake relasi~$\subseteq$
ing sawijining formula, kita bisa nggantine
nganggo formula kasebut (kanthi $x$ lan $y$ diganti
kaya perlune, lan variabel kaiket~$z$ diwenehi
jeneng anyar yen perlu). Upamane, ekstensionalitas
himpunan ateges yen sembarang himpunan~$x$ lan $y$
padha-padha ngemot sijine, mula $x$ lan $y$ iku
himpunan kang padha. Iki bisa dinyatakake kanthi
$\lforall[x][\lforall[y][((x \subseteq y \land y
    \subseteq x) \lif x = y)]]$, utawa, yen
$\subseteq$ diganti nganggo definisi ing ndhuwur,
kanthi
\[
\lforall[x][\lforall[y][((\lforall[z][(z \in x \lif z \in y)] \land
    \lforall[z][(z \in y \lif z \in x)]) \lif x = y)]].
\]
Iki pancen salah siji aksioma $\Log{ZFC}$, yaiku
"aksioma ekstensionalitas".

Ora ana !!{constant} kanggo $\emptyset$, nanging
kita bisa nyatakake "$x$ kosong" kanthi
$\lnot \lexists[y][y \in x]$. Banjur
"$\emptyset$ ana" dadi !!{sentence}~
$\lexists[x][\lnot \lexists[y][y \in
    x]]$. Iki aksioma liyane saka~$\Log{ZFC}$.
(Gatekna yen aksioma ekstensionalitas ngakibatake
yen mung ana siji himpunan kosong.) Kapan wae
kita arep ngrembug $\emptyset$ ing basa teori
himpunan, kita bakal nulis minangka "ana
himpunan kang kosong lan \dots" Upamane,
kanggo nyatakake yen $\emptyset$ iku
subhimpunan saka saben himpunan, kita bisa nulis
\[
\lexists[x][(\lnot\lexists[y][y \in x] \land \lforall[z][x \subseteq
    z])]
\]
ing ngendi, mesthi, $x \subseteq z$ kudu diganti
karo definisine.

Kanggo ngrembug operasi ing himpunan, kayata
$X \cup Y$ lan $\Pow{X}$, kita kudu nggunakake
cara kang padha. Ora ana simbol fungsi ing basa
teori himpunan, nanging kita bisa nyatakake relasi
fungsional $X \cup Y = Z$ lan $\Pow{X} = Y$ kanthi
\begin{align*}
& \lforall[u][((u \in x \lor u \in y) \liff u \in z)]\\
& \lforall[u][(u \subseteq x \liff u \in y )]
\end{align*}
amarga !!{element}s saka $X \cup Y$ persis
himpunan-himpunan kang dadi !!{element}s saka~$X$
utawa !!{element}s saka~$Y$, lan
!!{element}s saka $\Pow{X}$ persis subhimpunan
saka~$X$. Nanging iki ora marakake kita bisa
nggunakake $x \cup y$ utawa $\Pow{x}$ kaya-kaya
iku term: kita mung bisa nggunakake sakabehing
!!{formula}s kang nemtokake relasi
$X \cup Y = Z$ lan $\Pow{X} = Y$. Nyatane,
kita durung ngerti apa relasi-relasi iku tau
dicukupi, yaiku kita durung ngerti apa gabungan
lan himpunan pangkat tansah ana. Upamane,
!!{sentence} $\lforall[x][\lexists[y][\Pow{x} = y]]$
iku aksioma liyane saka~$\Log{ZFC}$
(aksioma himpunan pangkat).
Cathetan notasi: rumus kang nganggo tandha himpunan
pangkat ing kene mung cekakan metabahasa; ing formula
basa murni, relasi kasebut kudu dibentangake miturut
definisi ing ndhuwur.

Saiki kepriye anggone ngrembug pasangan urut
utawa fungsi? Ing kene kita kudu njlentrehake
kepriye pasangan urut lan fungsi bisa dianggep
minangka jinis himpunan tartamtu. Salah siji cara
nemtokake pasangan urut $\tuple{x, y}$ yaiku
minangka himpunan $\{\{x\}, \{x, y\}\}$.
Nanging kaya sadurunge, kita ora bisa ngenalake
!!a{function} kang menehi jeneng marang
himpunan iki; kita mung bisa nemtokake relasi
$\tuple{x, y} = z$, yaiku,
$\{\{x\}, \{x, y\}\} = z$:
\[
\lforall[u][(u \in z \liff (\lforall[v][(v \in u \liff v = x)] \lor
  \lforall[v][(v \in u \liff (v = x \lor v = y))]))]
\]
Iki nyatakake yen !!{element}s~$u$ saka~$z$
persis himpunan-himpunan kang mung nduweni $x$
minangka !!{element}, utawa mung nduweni $x$
lan~$y$ minangka !!{element}s (kanthi tembung liya,
himpunan kang padha karo $\{x\}$ utawa
padha karo $\{x, y\}$). Sawisé iku, kita bisa
nyatakake bab liyane, upamane, yen $X \times Y = Z$:
\[
\lforall[z][(z \in Z \liff \lexists[x][\lexists[y][(x \in
        X \land y \in Y \land \tuple{x, y} = z)]])]
\]

Fungsi $f \colon X \to Y$ bisa dianggep minangka
relasi $f(x) = y$, yaiku, minangka himpunan pasangan~
$\Setabs{\tuple{x,y}}{f(x) = y}$. Banjur kita bisa
kandha yen himpunan~$f$ iku fungsi saka $X$ menyang
$Y$ yen (a) iku relasi
$\subseteq X \times Y$, (b) total, yaiku
kanggo saben $x \in X$ ana $y \in Y$ saéngga
$\tuple{x, y} \in f$, lan (c) fungsional,
yaiku kapan wae $\tuple{x, y}, \tuple{x, y'} \in f$,
$y = y'$ (amarga nilai fungsi kudu tunggal).
Dadi " $f$ iku fungsi saka $X$ menyang $Y$ "
bisa ditulis minangka:
\begin{align*}
 \lforall[u][(u \in f \lif {}] & \lexists[x][\lexists[y][(x \in X \land y \in
      Y \land \tuple{x, y} = u)]]) \land {}\\
 \lforall[x][(x \in X \lif {}] &
      (\lexists[y][(y \in Y \land \mathrm{maps}(f, x, y))] \land {}\\
& (\lforall[y][\lforall[y'][((\mathrm{maps}(f, x, y) \land
    \mathrm{maps}(f, x, y')) \lif y = y')]]))
\end{align*}
ing ngendi $\mathrm{maps}(f, x, y)$ nyepetake
$\lexists[v][(v \in f \land
  \tuple{x, y} = v)]$ (!!{formula} iki
nyatakake "$f(x) = y$").

Saiki uga ora angel nyatakake yen
$f\colon X \to Y$ iku !!{injective}, upamane:
\begin{multline*}
 f \colon X \to Y \land \lforall[x][\lforall[x'][((x \in X \land x' \in
     X \land {}]] \\
 \lexists[y][(\mathrm{maps}(f, x, y) \land \mathrm{maps}(f,
        x', y))]) \lif x = x')
\end{multline*}
Fungsi~$f\colon X \to Y$ iku !!{injective}
yen lan mung yen, kapan wae $f$ masangake
$x, x' \in X$ karo siji~$y$, mula $x = x'$.
Yen formula iki dicekak minangka
$\mathrm{inj}(f, X, Y)$, kita wis bisa
nyatakake ing basa teori himpunan bab kang
ora prasaja kaya Teorema Cantor:
ora ana fungsi !!{injective} saka $\Pow{X}$
menyang $X$:
\[
\lforall[X][\lforall[Y][(\Pow{X} = Y \lif
    \lnot\lexists[f][\mathrm{inj}(f, Y, X)])]]
\]

Ana kang bisa ngira yen teori himpunan mbutuhake
aksioma liyane kang njamin anane himpunan kanggo
saben sipat kang nemtokake. Yen $!A(x)$ iku formula
teori himpunan kang variabel~$x$ katon bebas,
kita bisa nggatekake !!{sentence}
\[
\lexists[y][\lforall[x][(x \in y \liff !A(x))]].
\]
!!{sentence} iki nyatakake yen ana himpunan~$y$
kang !!{element}s-e persis kabeh $x$ kang
nyukupi~$!A(x)$. Skema iki diarani
"prinsip komprehensi". Katon migunani banget;
nanging sayange ora konsisten. Jupuken
$!A(x) \ident \lnot x \in x$, mula prinsip
komprehensi nyatakake
\[
\lexists[y][\lforall[x][(x \in y \liff x \notin x)]],
\]
yaiku nyatakake anane himpunan kabeh himpunan
kang dudu !!{element}s saka awake dhewe.
Himpunan kaya mangkono ora bisa ana---
iki Paradoks Russell. Nyatane, $\Log{ZFC}$
ngemot versi prinsip iki kang diwatesi---
ora njalari kontradiksi Russell iki---yaiku prinsip pemisahan:
\[
\lforall[z][\lexists[y][\lforall[x][(x \in y \liff (x \in z \land
      !A(x)))]]].
\]
Cathetan koreksi sumber SOL6-F010: kurung panutup matriks pepadhan kang ilang ing rumus pemisahan dibalekake. Cathetan cakupan SOL6-F011: klaim konsistensi sumber ora dibuktekake ing paragraf iki; pemisahan mung njupuk subset saka himpunan kang wis ana, mula instansi Russell iki ora njamin anane himpunan universal kang mbantah awake dhewe.

\begin{prob}
Tuduhna yen prinsip komprehensi ora konsisten
kanthi menehi !!a{derivation} kang nuduhake
\[
\lexists[y][\lforall[x][(x \in y \liff x \notin x)]] \Proves \lfalse.
\]
Mbokmenawa migunani yen luwih dhisik nuduhake
$(A \lif \lnot A) \land (\lnot A \lif A)
\Proves \lfalse$.
\end{prob}
\end{document}
